\section{Qualitätsanforderungen}

% Fertigen Sie hier eine Liste aller nicht-funktionalen Anforderungen an.
\begin{enumerate}
    \item Der Code soll lesbar und gut modifizierbar sein
    \item Die Ähnlichkeit wird in der Priorität: Titel, Graphiken, Bilder, dann Text bestimmt
    \item Es soll über VisualBasic auf die XML-Struktur der Präsentationen zugegriffen werden
\end{enumerate}

\subsection{Qualitätsziele des Projekts}

% Beschreiben Sie hier die einzelnen Qualitätsziele, die an die Software gestellt werden und wie diese am Ende des Projektes gemessen werden können. Beispielsweise könnte Sicherheit ein Qualitätsziel sein. Konkretiesieren Sie dieses und beschreiben Sie, was (hier Sicherheit) in Ihrer Software bedeutet und an welchen Stellen Sie besonders Acht geben müssen. Um zu beweisen, dass die Qualitätsziele erreicht wurden, müssen diese gemessen werden können. Überlegen Sie sich ein Verfahren, um die einzelnen Ziele messen zu können (bei Sicherheit könnte dies Beispielsweise ein Pentest sein.)

% Hier sind nur grobe Überlegungen notwendig. Tierfergehend werden Sie dies erst in SWQ kennenlernen.

\subsection{Qualitäts-Prioritäten des Kunden}

% Anhand der Informationen aus dem Video sollen die Anforderungen hier priorisiert werden. Wenn etwas unklar ist, fragen Sie Ihren Customer-Proxy

Die Qualitätsziele sind wie folgt absteigend priorisiert:
\begin{itemize}
    \item ...
\end{itemize}

\subsection{Wie Qualitätsziele erreicht werden sollen}
